% DRAFT branching-tree explainer for ex2 (Fig 1) — quick iteration.
% CHAIN layout: node contents defined as macros (top). Four chain HEADS are
% placed absolutely (e1, m1, t1, s1); every other node hangs off its parent
% at a fixed (dx, dy) offset — move a head and its whole chain follows; tune
% one offset to move one node (descendants follow). root/n1 hang off e1;
% l1 off m1; acl/js off ea. Edge labels attached to their \draw commands.
% Flush-left on purpose.
%
% STRUCTURAL CHANGE (flag if unwanted): e2 and a1 held identical configs, so
% they are MERGED into one node (ea) with two incoming edges — eager and lazy
% paths literally converge. Likewise the C# leaf merged into the
% Asyncio/Tokio leaf (same rule: [Block] destroys pending) => leaf (acl).
% The graph is now a DAG, not a tree.
\newcommand{\fgap}{\noalign{\vskip 3pt}}
% dimension highlights for rule names (colorbox: nestable in tikz nodes;
% colors match \ceag/\cext/\cdes from the prelude)
\newcommand{\dimeag}[1]{{\setlength{\fboxsep}{0.8pt}\colorbox{P1C1}{\textcolor{P1C1t}{#1}}}}
\newcommand{\dimext}[1]{{\setlength{\fboxsep}{0.8pt}\colorbox{P1C2}{\textcolor{P1C2t}{#1}}}}
\newcommand{\dimdes}[1]{{\setlength{\fboxsep}{0.8pt}\colorbox{P1C6}{\textcolor{P1C6t}{#1}}}}
\newcommand{\dimpro}[1]{{\setlength{\fboxsep}{0.8pt}\colorbox{P1C5}{\textcolor{P1C5t}{#1}}}}
\newcommand{\dimawa}[1]{{\setlength{\fboxsep}{0.8pt}\colorbox{P1C4}{\textcolor{P1C4t}{#1}}}}
\newcommand{\dimdir}[1]{{\setlength{\fboxsep}{0.8pt}\colorbox{P1C3}{\textcolor{P1C3t}{#1}}}}

%% ---------------- node contents (no layout) ----------------
%% FORMAL syntax (Section 4): \mathbfsf keywords, application by
%% juxtaposition, sys_io(t,v) for sleep, sys_block for the entry.
%% Named definitions keep the root un-inlined (work, process_detach, ex2).
\newcommand{\cfgRoot}{\begin{tabular}{@{}l@{}}
  $\mathit{work} \triangleq \mathbfsf{async}\ \mathbfsf{fn}\ ()\ \to$\\
  $\quad \mathbfsf{print}\ \text{"A"};\\\quad\mathbfsf{await}\ \mathbfsf{sys\text{-}io}(2, ());\\\quad\mathbfsf{print}\ \text{"B"}$\\
  \fgap
  $\mathit{process\text{-}detach} \triangleq \mathbfsf{async}\ \mathbfsf{fn}\ ()\ \to\\\quad\mathbfsf{spawn}\ (\mathit{work}())$\\
  \fgap
  $\mathit{ex2} \triangleq \mathbfsf{async}\ \mathbfsf{fn}\ ()\ \to$\\
  $\quad \mathbfsf{await}\ (\mathit{process\text{-}detach}());\\\quad \mathbfsf{await}\ \mathbfsf{sys\text{-}io}(1, ());\\\quad \mathbfsf{print}\ \text{"C"}$\\
  \fgap
  $\mathbfsf{sys\text{-}block}(\mathit{ex2}())$\\
  \fgap
  $\mathit{out}$: $\varepsilon$
  \end{tabular}}

% at work's application, process_detach inlined in main's frame
\newcommand{\cfgSpawnRedex}{\begin{tabular}{@{}l@{}}
  $C[\, \mathbfsf{spawn}\ (\mathit{work}()) \,]$\\
  \fgap
  $\mathit{out}$: $\varepsilon$
  \end{tabular}}

% work queued, unrun; plain C = no scope-exit obligation (indefinite extent)
\newcommand{\cfgQueued}{\begin{tabular}{@{}l@{}}
  $C[\, t_{\mathit{work}} \,]$\\
  \fgap
  $Q$: $\mathbfsf{print}\ \text{"A"};$\\
  \hphantom{$Q$: }$\mathbfsf{await}\ \mathbfsf{sys\text{-}io}(2, ());$\\
  \hphantom{$Q$: }$\mathbfsf{print}\ \text{"B"}$\\
  \fgap
  $\mathit{out}$: $\varepsilon$
  \end{tabular}}

% work queued with the INSTALLED scope-exit epilogue visible inline:
% green await = extent obligation; coral cancel (Swift) = destruction
\newcommand{\cfgQueuedTrio}{\begin{tabular}{@{}l@{}}
  $C[\, t_{\mathit{work}};$\\
  \hphantom{$C[\,$}\dimext{$\mathbfsf{await}\ t_{\mathit{work}}$}$\,]$\\
  \fgap
  $Q$: $\mathbfsf{print}\ \text{"A"};$\\
  \hphantom{$Q$: }$\mathbfsf{await}\ \mathbfsf{sys\text{-}io}(2, ());$\\
  \hphantom{$Q$: }$\mathbfsf{print}\ \text{"B"}$\\
  \fgap
  $\mathit{out}$: $\varepsilon$
  \end{tabular}}

\newcommand{\cfgQueuedSwift}{\begin{tabular}{@{}l@{}}
  $C[\, t_{\mathit{work}};$\\
  \hphantom{$C[\,$}\dimdes{$\mathbfsf{cancel}\ t_{\mathit{work}}$}$;$\\
  \hphantom{$C[\,$}\dimpro{$\mathbfsf{try}$}\\
  \hphantom{$C[\,$}\quad\dimext{$\mathbfsf{await}\ t_{\mathit{work}}$} \\
  \hphantom{$C[\,$}\dimpro{$\mathbfsf{catch\ fn}\ \mathit{\text{-}}\ \to\ ()$}
  $\,]$\\
  \fgap
  $Q$: $\mathbfsf{print}\ \text{"A"};$\\
  \hphantom{$Q$: }$\mathbfsf{await}\ \mathbfsf{sys\text{-}io}(2, ());$\\
  \hphantom{$Q$: }$\mathbfsf{print}\ \text{"B"}$\\
  \fgap
  $\mathit{out}$: $\varepsilon$
  \end{tabular}}

% lazy intermediate: coroutine made ([Async-App]), spawn is the redex;
% the Spawn VARIANT (dyn vs indef) decides the next branch
\newcommand{\cfgLazyCor}{\begin{tabular}{@{}l@{}}
  $C[\, \mathbfsf{spawn}\ \mathit{coro}_{\mathit{work}} \,]$\\
  \fgap
  $\mathit{out}$: $\varepsilon$
  \end{tabular}}

% Swift abort snapshot: work ran (A printed), its cancellation-aware sys_io
% observed the flag and raises; main still blocked on the epilogue await
\newcommand{\cfgSwiftAbort}{\begin{tabular}{@{}l@{}}
    $C$$[\,\dimawa{\mathbfsf{throw}\ \text{"cancelled"}}\,]$\\
  \fgap
  $\mathit{out}$: A
  \end{tabular}}

% eager: work already ran to its await; suspended in T
\newcommand{\cfgEagerT}{\begin{tabular}{@{}l@{}}
  $C[\, t_{\mathit{work}} \,]$\\
  \fgap
  $T$: $(2,\ \mathit{work},\ \mathbfsf{print}\ \text{"B"})$\\
  \fgap
  $\mathit{out}$: A
  \end{tabular}}

% main done; orphan pending in T; sys_block holds the result
\newcommand{\cfgBlockDone}{\begin{tabular}{@{}l@{}}
  $\mathbfsf{sys\text{-}block}(())$\\
  \fgap
  $T$: $(2,\ \mathit{work},\ \mathbfsf{print}\ \text{"B"})$\\
  \fgap
  $\mathit{out}$: AC
  \end{tabular}}

% work settled/gone; only main's tail remains (out varies). No more C:
% the remaining code is spelled out (the context here is empty)
\newcommand{\cfgSettled}[1]{\begin{tabular}{@{}l@{}}
  $();$\\
  $\mathbfsf{await}\ \mathbfsf{sys\text{-}io}(1, ());$\\
  $\mathbfsf{print}\ \text{"C"}$\\
  \fgap
  $\mathit{out}$: #1
  \end{tabular}}

% leaves: output + runtimes
\newcommand{\tleaf}[2]{\begin{tabular}{@{}c@{}}#1\\{\rmfamily #2}\end{tabular}}

%% ---------------- layout ----------------
% NOTE: flush-left on purpose during iteration — a centered figure re-balances
% whenever any node shifts, moving everything else. Re-center at the very end
% (wrap in center, or \useasboundingbox to freeze the bbox).
\noindent
% {\tiny where $C \triangleq [\cdot];\ \mathbfsf{await}\ \mathbfsf{sys\text{-}io}(1, 0);\ \mathbfsf{print}\ \text{"C"}$ --- the rest of $\mathit{ex2}$; $\mathbfsf{sys\text{-}block}$ is elided until \semrule{Block-Done} applies}\par\smallskip

\begin{tikzpicture}[
  cfg/.style={draw=black!20, rounded corners=2pt, inner sep=2.5pt, font=\tiny\ttfamily},
  leaf/.style={draw=black!40, thick, rounded corners=2pt, inner sep=2.5pt, font=\tiny\ttfamily},
  lab/.style={font=\tiny, align=left, inner sep=1pt, fill=white, fill opacity=0.85, text opacity=1},
  labr/.style={lab, align=right},
  arr/.style={->, >={Stealth[length=1.4mm]}, semithick}
]
  %% Chain heads (absolute): move a head and its whole chain follows.
  \node[cfg, anchor=north] (e1) at (1.4, 0)      {\cfgEagerT};    % eager lane
  \node[cfg, anchor=north] (m1) at (4.95, -1.65) {\cfgSettled{$\varepsilon$}}; % smol lane
  \node[cfg, anchor=north] (lt) at (7.55, 1.20) {\cfgQueuedTrio}; % trio lane
  \node[cfg, anchor=north] (s1) at (10.85, 2.5)  {\cfgQueuedSwift}; % swift lane

  %% Chains: each node hangs off its parent's north at a fixed (dx, dy);
  %% tune the offset to move one node (its own descendants follow).
  \node[cfg,  anchor=north] (ea) at ([shift={(0,-2.15)}]e1.north) {\cfgBlockDone};
  \node[leaf, anchor=north] (sm) at ([shift={(0,-2.4)}]m1.north) {\tleaf{C}{Smol}};
  \node[cfg,  anchor=north] (t2) at ([shift={(0,-3.0)}]lt.north) {\cfgSettled{AB}};
  \node[leaf, anchor=north] (tr) at ([shift={(0,-2.4)}]t2.north)  {\tleaf{ABC}{Trio}};
  \node[cfg,  anchor=north] (sab) at ([shift={(0,-3.4)}]s1.north) {\cfgSwiftAbort};
  \node[cfg,  anchor=north] (s3) at ([shift={(0,-1.6)}]sab.north) {\cfgSettled{A}};
  \node[leaf, anchor=north] (sw) at ([shift={(0,-2.4)}]s3.north)  {\tleaf{AC}{Swift}};


  % ea's two leaves sit side by side under it (outside the grid), so both
  % edges can fan from the same point without crossing a box
  \node[leaf, below=6mm of ea.south west, anchor=north] (acl) {\tleaf{AC}{\CSharp, Asyncio, Tokio}};
  \node[leaf, below=10mm of ea.south east, anchor=north] (js) {\tleaf{ACB}{\Js}};

  % the two wide nodes live outside the grid, above it
  \node[cfg, above=12mm of e1.north west, anchor=south west] (root) {\cfgRoot};
  \node[cfg, right=22mm of root.north east, anchor=north west, yshift=-6mm] (n1) {\cfgSpawnRedex};
  % directly below n1; the one number = border-to-border gap
  \node[cfg, below=9mm of n1] (l0) {\cfgLazyCor};
  % smol/asyncio/tokio: plain-C queued state, below-left of l0
  \node[cfg, anchor=north] (l1) at ([shift={(-2.0,-0.5)}]l0.south) {\cfgQueued};

  %% ---------------- edges (all straight; labels attached) ----------------
  % start on root's edge at n1's height => perfectly horizontal
  \draw[arr] (root.east |- n1.west) -- (n1.west)
    node[lab, midway, below=2pt] {\semrule{Block-Wait}\\ \semrule{Async-App}\\ \semrule{Await}};

  \draw[arr] (n1.south) -- (s1.north)
    node[labr, pos=0.6, midway, right=2pt] {\dimeag{\semrule{Async-App}}\textsubscript{\designpoint{semi-eager}}};

  \draw[arr] (n1.south) -- (l0.north)
    node[labr, near end, right=2pt] {\dimeag{\semrule{Async-App}}\textsubscript{\designpoint{lazy}}};

  % Spawn variants: dyn installs the scope-exit obligation (C turns green)
  \draw[arr] (l0.south) -- (lt.north)
    node[lab, midway, right=2pt] {\semrule{Spawn}};
  \draw[arr] (l0.south) -- (l1.north)
    node[labr, midway, left=2pt] {\semrule{Spawn}};

  \draw[arr] (n1.south) -- (e1.north)
    node[labr, pos=0.6, near end, left=2pt, align=left] {\dimeag{\semrule{Async-App}}\textsubscript{\designpoint{eager}}\\ \semrule{OS-IO}};

  \draw[arr] (e1.south) -- (ea.north)
    node[lab, midway, right=2pt] {\semrule{OS-IO}\\ \semrule{Signal}\\ \semrule{Schedule}};

  \draw[arr] (s1.south) -- (sab.north)
    node[labr, midway, right=2pt, align=left] {\semrule{Schedule}\\ \semrule{OS-IO}};

  \draw[arr] (sab.south) -- (s3.north)
    node[labr, midway, right=2pt, align=left] {\semrule{Await}\\ \semrule{Catch-Exn}};

  \draw[arr] (s3.south) -- (sw.north)
    node[labr, midway, right=2pt, align=left] {\semrule{OS-IO}\\ \semrule{Signal}\\ \semrule{Schedule}\\ \semrule{Block-Done}};

  \draw[arr] (l1.south) -- (m1.north)
    node[labr, near end] {\dimdes{\semrule{Cancel-Unstarted}}};

  \draw[arr] (m1.south) -- (sm.north)
    node[lab, midway, right=2pt] {\semrule{OS-IO}\\ \semrule{Signal}\\ \semrule{Schedule}\\ \semrule{Block-Done}};

  \draw[arr] (lt.south) -- (t2.north)
    node[lab, midway, right=2pt] {\semrule{Schedule}\\ \semrule{OS-IO}\\ \semrule{Signal}\\ \semrule{Await-Task}};

  \draw[arr] (t2.south) -- (tr.north)
    node[lab, midway, right=2pt] {\semrule{OS-IO}\\ \semrule{Signal}\\ \semrule{Schedule}\\ \semrule{Block-Done}};

  \draw[arr] (l1.south) -- (ea.north east)
    node[lab, pos=0.4, near end, above=2pt] {\semrule{Schedule}\\ \semrule{OS-IO}\\ \semrule{Signal}};

  \draw[arr] (ea.south) -- (acl.north)
    node[lab, midway, left=2pt] {\dimdes{\semrule{Block-Done}}\textsubscript{\designpoint{terminated}}};

  \draw[arr] (ea.south) -- (js.north)
    node[lab, midway, right=2pt] {\semrule{Signal}\\ \semrule{Schedule}\\ \dimdes{\semrule{Block-Done}}\textsubscript{\designpoint{awaited}}};
\end{tikzpicture}
